% !TeX program = xelatex
% 省墨讲义模板：单栏、5 号字、1.5cm 页边距、细线框、无彩色填充。
% 适合打印；内容连续排版，章节不强制换页。
\documentclass[a4paper,oneside,zihao=5,fontset=fandol]{ctexbook}

% ===== 标题 =====
\newcommand{\handouttitle}{《课程名称》复习精要}
\newcommand{\handoutsubtitle}{知识点汇总 \quad 适用：期末复习}

% ===== 版式 =====
\usepackage{amsmath,amssymb}
\usepackage[margin=1.5cm,headheight=1cm,footskip=0.8cm]{geometry}
\usepackage{enumitem}
\usepackage{titlesec}
\usepackage{tcolorbox}
\tcbuselibrary{skins,breakable}
\usepackage{multicol}   % 需要双栏速查时使用

\linespread{1.2}
\setlist{nosep}
\pagestyle{plain}
\raggedbottom

% 章节连续排版（省纸）
\ctexset{chapter/break=\par}

\titleformat{\section}
  {\normalfont\large\bfseries}
  {\thesection}{1em}{}[\titlerule]
\titleformat{\subsection}
  {\normalfont\normalsize\bfseries}
  {\thesubsection}{0.5em}{}

% 细线重点框：不填色，只画线
\newtcolorbox{mybox}[1]{
  enhanced, breakable,
  colback=white, colframe=black!55,
  colbacktitle=white, coltitle=black,
  boxrule=0.3pt, leftrule=2pt, arc=0mm,
  fonttitle=\bfseries\small, title={#1},
  left=4pt, right=4pt, top=3pt, bottom=3pt,
  before skip=6pt, after skip=6pt
}

\begin{document}

% ===== 首页 =====
\vspace*{0.8cm}
\begin{center}
  {\Huge\bfseries \handouttitle\par}
  \vspace{0.3cm}
  {\large \handoutsubtitle\par}
\end{center}
\vspace{0.6cm}

\chapter{基础概念}

\begin{mybox}{本章重点}
用两三句话概括本章必须掌握的概念、公式与易错点。
\end{mybox}

\section{概念与特征}

先写定义和适用范围，再列关键特征。用 \textbf{粗体} 标出术语或结论。

\begin{itemize}[left=0pt]
  \item \textbf{定义}：术语的含义与使用场景。
  \item \textbf{特征}：判断与区分的条件。
  \item \textbf{例子}：一个简短的应用情境。
\end{itemize}

\subsection{概念辨析}

对容易混淆的术语，分别记录适用范围、典型例子和判断依据。

\section{步骤与方法}

\begin{enumerate}[left=0pt]
  \item \textbf{明确对象}：确定问题、范围和已知条件。
  \item \textbf{选择方法}：说明方法与适用条件。
  \item \textbf{检查结果}：把结论与题目要求逐项对照。
\end{enumerate}

\chapter{应用与练习}

\section{典型题型}

先用一句话概括题目要求，再整理答题要点。

\subsection{简答题}

\begin{itemize}[left=0pt]
  \item \textbf{答题顺序}：定义 → 分点解释 → 例子。
  \item \textbf{检查项目}：术语准确，每个要点都回应题意。
\end{itemize}

\end{document}
